\chapter{Theoretical Framework and Related Work}
\label{ch:literature}

This chapter situates the thesis within four bodies of work: the Node--Place tradition
that frames transit stations as points where a transport network and an urban place
coincide; travel-demand estimation, which supplies the demand signal the framework is
built on; accessibility and coverage-gap analysis, which turns demand and supply into a
diagnostic; and the machine-learning and network-design literature from which the
predictive and generative components are drawn. The chapter closes by locating the
research gap---a demand-driven, non-circular, transferable pipeline---that the remainder
of the thesis addresses.

\section{The Node--Place model and its extensions}
\label{sec:lit-nodeplace}

The conceptual anchor of this work is Bertolini's Node--Place model
\parencite{bertolini1996, bertolini1999}, which characterizes a station area along two
axes: its \emph{node} value, the position it occupies in the transport network, and its
\emph{place} value, the intensity and diversity of the urban activities around it. The
model's central claim is that sustainable transit-oriented development requires the two
to stay in balance; persistent imbalance signals either an under-used station in a rich
urban setting or an over-provided node in a weak one, a reading later operationalized as
an empirical station classification \parencite{reusser2008, vale2015}. The place axis
itself draws on a large built-environment literature relating density, land-use
diversity, and design to travel behavior \parencite{cervero1997threeD, ewing2010meta},
which motivates the place indicators used here. As a diagnostic the model is attractive
because it is descriptive and low-data---node and place indices can be assembled from
network and land-use variables alone---but in its original form it is static and says
nothing about \emph{latent demand} or about where new nodes ought to be placed.

Recent extensions add dimensions to the original two, generally coupling the enlarged
indicator set with machine learning. \textcite{su2022nodeplace} augment node and place
with a ridership term and interpret it with \textsc{shap}, moving the tradition from
description toward prediction; \textcite{aminipishro2022nprt} instead add both ridership
and time-of-day as a fourth and third dimension (an NP--RT rather than NP--RV
formulation), classifying stations with a K-means/Cube method. Neither adds a
\emph{vitality} term specifically, and no verified paper matching an ``NP--RV''
node--place--ridership--vitality formulation could be located during a 2026-09-09
literature-review pass; an earlier draft of this chapter cited such a paper, which is
removed here rather than left unverifiable. The broader pattern the two verified papers
above illustrate---augmenting Node--Place with an outcome variable and an interpretable
learner---is the one this thesis follows structurally, while departing from it on the
outcome variable itself (\cref{sec:lit-ml}).

Bertolini's own later extension pairs the model with an explicit accessibility measure
across multiple European metropolitan areas rather than a ridership/vitality term
\parencite{papa2015tod}, which is the closer lineage antecedent to the accessibility- and
coverage-gap-based extension developed in this thesis (\cref{sec:w3}). This thesis adopts
a Node--Place--People (\nppv) variant: it retains node and place, foregrounds a
\emph{people} dimension (population, dependency, and equity attributes), and---crucially---
removes the vitality term from the prioritization weighting because, in the study area, the
only available vitality proxy is resolved at the municipality level and therefore collapses
to a binary ``has-premium-transit'' flag that reintroduces exactly the supply-following
circularity the thesis sets out to avoid (\cref{sec:w4}). This exclusion is reported here as
a data-availability limitation of the specific vitality proxy available in \zmg, not as a
theoretical claim that vitality is dispensable in principle for a Node--Place--type model;
an AGEB-resolution vitality proxy (e.g., from pedestrian counts or a demand-independent
activity-intensity measure) remains an open direction, and the renamed \nppv{} label should
be read as scoped to this application rather than as a general claim that a vitality or
ridership dimension is dispensable for Node--Place extensions generally.

\section{Travel demand estimation}
\label{sec:lit-demand}

The demand layer follows the classical four-step tradition
\parencite{mcnally2000fourstep, ortuzar2011}, of which the two steps used here are trip
generation and trip distribution. Trip distribution is modeled with a spatial-interaction
(gravity) model, whose modern statistical justification is the entropy-maximizing
derivation of \textcite{wilson1971gravity}: among all origin--destination matrices
consistent with the observed marginal totals, the gravity form is the most probable, and
its distance-decay parameter governs how sharply interaction falls with separation. The
doubly-constrained variant, balanced by the iterative proportional-fitting (Furness)
procedure \parencite{furness1965}, reproduces both origin and destination totals exactly
and is the form adopted in \cref{sec:w1}.

Two design choices distinguish the demand layer here from a textbook application. First,
trip \emph{generation} is built entirely from Tier-1 census and economic variables, with
no transit input, so that demand is never a function of current supply---the property
the whole framework depends on. Second, the distance-decay parameter is not assumed but
\emph{calibrated} against an independent metropolitan travel survey
\parencite{imeplan_eod2022}, so that the impedance function reflects observed behavior
rather than a conventional exponent (\cref{sec:w2}). Calibration also yields a documented
transfer error, which matters when the same prior is carried to other cities
(\cref{ch:transferability}).

\section{Accessibility and coverage-gap analysis}
\label{sec:lit-accessibility}

Supply is measured with a cumulative-opportunities accessibility surface: for each unit,
the number of opportunities (here, jobs) reachable within a fixed travel-time budget over
the existing transit network. Accessibility---the ease of reaching valued
destinations---has a long conceptual and methodological literature
\parencite{handy1997, geurs2004}; cumulative-opportunities measures are among its most
widely used forms because they are transparent and require only a network and an
opportunity distribution, at the cost of a hard cutoff at the budget boundary. The
travel-time budget here combines access walk time, expected wait derived from headways,
and in-vehicle time (\cref{sec:w3}).

The diagnostic contribution of the thesis lives at the intersection of demand and
accessibility. Rather than label areas by whether they currently have service---the
circular target criticized in \cref{sec:intro-circularity}---the framework defines a
\emph{coverage gap} as modeled demand set against measured accessibility, so that a unit
scores high only when latent demand is genuinely unmet by current supply. This connects
to two related strands: the ``transit desert'' literature, which contrasts transit supply
against demand at fine spatial resolution \parencite{jiao2013deserts}, and the
social-need gap tradition, which weights under-supply by the social needs of the
population it affects \parencite{currie2010gaps}. Closer still is
\textcite{almamun2011servicegaps}'s accessibility-based transit need index, which compares
a need index against an accessibility index at a similarly fine resolution specifically to
identify service gaps---the same two-index-comparison structure this thesis's coverage-gap
ratio (\cref{eq:gap}) implements. The framework's contribution relative to this existing
gap-index tradition is not the comparison structure itself, which is well established, but
feeding the resulting gap directly forward as a supervised-learning target
(\cref{sec:lit-ml}) rather than stopping at a descriptive map. Because demand and supply
are constructed independently, this produces a dependent variable that a supervised model
can learn without simply memorizing the existing network.

\section{Multi-criteria decision analysis and objective weighting}
\label{sec:lit-mcda}

Turning a set of indicators into a single priority score requires weights. Multi-criteria
decision analysis is well established in transport-project evaluation
\parencite{macharis2015}, but expert- or survey-elicited weights import analyst bias and
are hard to transfer between cities. This thesis instead uses objective weighting,
combining the CRITIC method \parencite{diakoulaki1995critic}---which weights a criterion
by its contrast intensity and its independence from the others---with the Entropy Weight
Method, which weights by the dispersion of each criterion's values as measured by Shannon
entropy \parencite{shannon1948}. Both are data-driven in the specific sense that the final
weight-assignment step reads the weights off the indicator matrix itself, which makes the
prioritization reproducible and re-derivable from scratch in a new city without
re-eliciting expert judgments \emph{for that step}. This does not remove analyst judgment
from the pipeline altogether: which fourteen indicators enter the matrix, how they are
grouped into the Node/Place/People taxonomy, which normalization rule (log-transform versus
plain min--max) each receives, and the equal-weighted CRITIC--EWM ensemble average are all
upstream choices made by the analyst, not read off the data. \Cref{sec:area-integrity}'s
own account of a marginalization-index sign-and-imputation correction, which shifted the
mean equity score by several tenths, is direct evidence of how consequential these upstream,
non-``objective'' choices can be; CRITIC/EWM removes expert judgment from the final
aggregation step, not from the framework as a whole. Equity is
handled not by manipulating these weights but by an explicit, separable additive term
built from marginalization and social-lag indices
\parencite{conapo_marginacion2020, coneval_rezago2020}, so that its contribution to the
final score is transparent and can be reported with and without (\cref{eq:equity}). This
separable treatment reflects a growing view that transport equity should be an explicit
objective rather than an implicit by-product \parencite{martens2016, pereira2017justice}.
The qualification above---that CRITIC/EWM's objectivity is confined to the final
aggregation step---is not merely a self-critique: \textcite{mukhametzyanov2021critic}'s
comparative analysis of Entropy, CRITIC, and standard-deviation weighting finds that
formal decision-matrix-based weighting is not always well-behaved, which is independent
confirmation that these methods' ``objective'' framing needs the same qualification this
thesis applies to it. This combination of CRITIC-family weighting with a separate
gap/accessibility analysis is not unique to this thesis: \textcite{rathod2024critic}
apply the same pairing---CRITIC-MCDM weighting alongside a transit accessibility
performance-gap analysis---to Surat, India, a 2024 precedent for treating the
prioritization layer and the gap diagnostic as related but distinct components, as this
thesis does across W3 and W4.

A specific, checkable failure mode the literature above identifies is high correlation
among indicators: CRITIC's weight formula discounts a criterion for correlating with
others, but the Entropy Weight Method does not, so an ensemble average of the two only
partially corrects for redundant indicators. A correlation check of the fourteen
Node--Place--People indicators (recomputed 2026-09-09) finds three pairs above the common
$|r|>0.8$ multicollinearity threshold---\code{p\_poi\_density\_n}~/~\code{p\_retail\_
density\_n} ($r=0.93$), \code{n\_intersections\_n}~/~\code{n\_street\_density\_n}
($r=0.828$), and \code{pe\_population\_n}~/~\code{pe\_pop\_density\_n} ($r=0.812$)---out
of a mean absolute off-diagonal correlation of \num{0.335}. None of the three pairs is
surprising (each pairs a count with a closely related density or rate), and CRITIC's own
weighting already discounts them; the check is reported here because it is the concrete
diagnostic the failure-mode literature calls for, not because it overturns the weighting
result.

\section{Machine learning for station and corridor siting}
\label{sec:lit-ml}

A substantial literature applies supervised learning to transit siting.
\textcite{niu2023rf} use random forests to predict station suitability, and the
Node-Place-machine-learning line reviewed above \parencite{su2022nodeplace,
aminipishro2022nprt} pairs tree-based models with \textsc{shap} values
\parencite{lundberg2017shap} for interpretability. The framework here uses the same tools
---random forests \parencite{breiman2001rf} and gradient-boosted trees
\parencite{chen2016xgboost, ke2017lightgbm}, read through \textsc{shap}---but
departs from much of this literature on the dependent variable. Where a model is trained
to predict a ``has-a-stop'' label, or a suitability class that is itself a function of
current supply, it can only learn to reproduce the status quo and will report high
accuracy for a circular reason. The thesis makes this failure mode explicit: an earlier
version of its own pipeline trained on a cluster-derived label and reported near-perfect
recovery accuracy that carried no planning information, and it was retired in favor of a
model trained on the independent coverage-gap target (\cref{sec:w3}).

This specific transit-siting circularity is not an isolated observation about one
pipeline's history; it is a special case of a mechanism given precise formal treatment in
three more mature literatures, introduced in \cref{sec:intro-circularity} and elaborated
here. \textcite{kaufman2012leakage} name and formalize the general data-mining failure
mode---\emph{leakage}, the use of target-derived information the model should not
legitimately have access to---of which a supply-derived siting label is one instance.
\textcite{lakkaraju2017selectivelabels}'s \emph{selective labels problem} gives the
sharpest statistical account of exactly why this particular instance of leakage is
unavoidable by construction rather than merely a data-quality defect: whenever an outcome
is observed only conditional on a prior decision (their bail-release setting; here, a
prior siting decision), no amount of additional data or modeling sophistication recovers
the counterfactual outcome for units where the decision went the other way. Empirically,
\textcite{lum2016predict} and, more rigorously, \textcite{ensign2018runaway} show that
predictive-policing models trained on where police were previously sent reproduce and can
amplify that historical pattern rather than detecting where crime genuinely concentrates,
a runaway feedback loop with the identical logical structure---train on where the
intervention already happened, and the model can only relearn the intervention's own
footprint---and \textcite{flynn2022spacetime} raise the same concern specifically for
spatial infrastructure allocation. Grounding the thesis's circularity argument in this
established, cross-disciplinary body of work, rather than treating it as an ad hoc
observation about one pipeline's history, is one of the corrections made during a
2026-09-09 literature-review pass. The methodological point---that interpretability tools
cannot rescue a selectively-labeled target, and that the target itself must be constructed
independently of the prior decision it would otherwise encode---is a central lesson of
this work, and one with precedent well beyond transit planning.

\section{Network design and corridor generation}
\label{sec:lit-network}

Generating candidate corridors from a set of priority locations is a network-design
problem, part of the broad literature on transit-network design
\parencite{guihaire2008, farahani2013review}, of which \textcite{kepaptsoglou2009review}'s
review---organizing the field by design objectives, operating-environment parameters, and
solution approach---remains the standard reference point for classifying a new method such
as this thesis's W5/W6 combination. The relevant heuristic tradition is the
Steiner tree in graphs---connecting a set of required (terminal) nodes at minimum cost,
possibly through additional intermediate nodes---for which \textcite{takahashi1980} give
a classic shortest-path-based approximation and \textcite{kmb1981steiner} a fast
minimum-spanning-tree-based one, the family of relatives used here. Spanning-tree-based
transit design is not purely historical: a 2026 case study applies the same underlying
idea---a tabu-search-refined spanning tree minimizing total passenger-kilometres---to the
Canberra bus network \parencite{suguira2026spanningtree}, evidence that the heuristic
family remains in active use for real-network (re)design, though as a preprint that
result should be weighted accordingly. These heuristics are attractive for corridor
generation because they operate directly on the road graph and produce connected,
buildable geometry from scattered demand anchors.

Their limitation, which the thesis characterizes in detail, is that a tree built to
minimize connection cost among anchors is not the same object as a good \emph{corridor}.
A spanning tree walked into a single line can insert phantom jumps between non-adjacent
branches, and---more consequentially---a feasibility filter based on endpoint detour will
select for few, near-collinear anchors regardless of the demand a corridor actually
serves, confounding feasibility with anchor sparsity. \Cref{sec:w6} re-architects the
generator around a diameter-trunk shaper and an anchor-directness feasibility metric that
measure the right thing, and \cref{ch:discussion} reports both what this fixes and the
residual limitation that remains. A separate positioning point, relative to
\textcite{kepaptsoglou2009review}'s taxonomy, is more favorable to this thesis's design
than an earlier draft of this section suggested. Academic multi-objective TNDP research
has moved toward many-objective Pareto exploration that avoids imposing fixed weights at
all, and \textcite{silva2022mo} note this directly: it is genuinely difficult for
transportation experts to act on a raw Pareto front, or to combine objectives in a way
that adapts flexibly as priorities change, which is precisely why applied practice still
leans on weighted-sum scalarization despite its inability to reach non-convex regions of
the trade-off space. Rather than treating one approach as simply superseding the other,
\textcite{silva2022mo} argue single- and multi-objective framings are best combined
synergistically. The composite score in \cref{eq:objectives} does exactly this: it blends
$f_1$, efficiency, and $f_3$ at fixed weights (0.50/0.25/0.25) \emph{alongside}, not
instead of, genuine Pareto ranking (\cref{fig:pareto}), giving a reader both an
actionable single number and the full non-dominated set to inspect. This hybrid is
consistent with, rather than behind, how the field currently reconciles academic and
applied practice on this question---though, as \cref{sec:res-w6,sec:disc-generative}
show, having both is only useful if the composite score and the Pareto ranking actually
agree, which for this thesis's own candidate set they do not.

\section{Research gap}
\label{sec:lit-gap}

The four literatures above are individually mature but are rarely combined into a single,
reproducible pipeline that is also \emph{non-circular} and \emph{transferable}. Demand
models estimate flows but are seldom fed forward into an explicit gap diagnostic; Node--Place
studies describe stations but do not generate corridors; machine-learning siting work
frequently trains on supply-derived targets; and corridor-generation heuristics are rarely
evaluated against genuine unmet need rather than against the network that already exists.
This is not to say the diagnostic half of this claim is without antecedent:
\textcite{foth2013equitable} construct essentially the same accessibility-versus-social-need
overlay used here as a spatial equity diagnostic for Toronto, and \textcite{camporeale2019equity}
directly address network-design solutions evaluated against horizontal and vertical equity
principles rather than cost or coverage alone---the same territory this thesis's generative
layer (\cref{sec:w5,sec:w6}) occupies. What is not combined in either of those works, nor
elsewhere in the reviewed literature, is the specific chain this thesis closes: a Tier-1-only
demand estimate, fed into a non-circular supervised retrain, coupled to an MST-diameter-trunk
generator gated on anchor-directness, and the whole chain transferred end-to-end across
metropolitan areas of contrasting size. The gap this thesis addresses is therefore narrower
than ``an integrative pipeline'' per se---closely related integrative diagnostics and
equity-aware generators already exist---and is better stated as: a demand-driven framework
that (i) estimates demand from Tier-1 data alone, (ii) diagnoses the coverage gap with a
target constructed independently of supply and validates it out-of-sample, (iii) generates
and evaluates candidate corridors on need, non-redundancy, and demand efficiency, and (iv)
transfers end-to-end to metropolitan areas of contrasting size without bespoke
re-engineering. The methodology in \cref{ch:methodology} is organized to close exactly this
narrower, but still unaddressed, combination.
